% Chapter 10: Evaluation Methodology

\chapter{Evaluation Methodology}

This chapter outlines an evaluation plan for PVE-ZKP. It defines the goals and
metrics of interest, describes an experimental setup for measuring performance
and resource usage, and sketches a set of experiments that can be carried out
once the implementation is deployed in a suitable environment.

\section{Evaluation Goals}

The primary goals of the evaluation are:

\begin{itemize}[leftmargin=*]
  \item to quantify the cost of proving and verification for the five policy
    circuits used in PVE-ZKP;
  \item to measure end-to-end latency for form submissions under realistic
    workloads;
  \item to assess how performance scales with increasing numbers of concurrent
    clients; and
  \item to compare PVE-ZKP against a baseline architecture that performs
    validation entirely on the receiving system without zero-knowledge proofs.
\end{itemize}

\section{Metrics}

Relevant metrics include:

\begin{itemize}[leftmargin=*]
  \item proving time per circuit and per submission;
  \item verification time per circuit and per submission;
  \item memory usage and CPU utilization on the validation service and
    receiving system;
  \item end-to-end response time as perceived by the client; and
  \item throughput (submissions per second) under varying loads.
\end{itemize}

These metrics can be collected using standard profiling and benchmarking tools
available in Node.js and the underlying operating system.

\section{Experimental Setup}

A typical experimental setup would fix:

\begin{itemize}[leftmargin=*]
  \item hardware characteristics (CPU, memory, storage);
  \item network conditions (local vs remote deployment, latency, bandwidth);
  \item versions of Circom, snarkjs, and other dependencies; and
  \item representative input distributions for the healthcare-style form.
\end{itemize}

Experiments should be automated via scripts that deploy the validation and
receiving services, generate synthetic but realistic workloads, and record
results for later analysis.

\section{Planned Experiments}

Planned experiments include:

\begin{itemize}[leftmargin=*]
  \item microbenchmarks that measure prove and verify times for each individual
    policy circuit across a range of input sizes and parameters;
  \item end-to-end benchmarks that measure client-perceived latency and
    throughput as the number of concurrent form submissions increases; and
  \item comparative benchmarks that implement equivalent policies directly on
    the receiving system without ZKPs to quantify overheads introduced by
    PVE-ZKP.
\end{itemize}

Results from these experiments would be incorporated into the thesis as tables
and figures once available, providing a quantitative complement to the
qualitative security analysis presented earlier.
